\scanpage{265}
\mspage{77}
% archive sheet number: 265
En fait, on voit \deleted{\unsure{par là}} que le foncteur

\[
i_{1*}i_{0!}:\quad\mathcal A\simeq\mathbb D(\deleted{\unsure{?}}I_0)\longrightarrow\mathbb D(I)\deleted{\unsure{\mathrm{sub}}}
\]

induit une équivalence de $\mathcal A$ sur la sous-catégorie pleine de $\mathbb D(I)$ formée des carrés $\mathbb D$-cartésiens \deleted{tels que} dont les fibres en $a$ et $b$ sont nulles. \inserted{(le foncteur quasi-inverse évident est $i^*$, où $i:I_0\lhook\joinrel\longrightarrow I$.)}

Cela signifie donc que dans ce cas, la donnée du squelette \inserted{de nature} très \unsure{particulier|particulière} ($F_a\simeq F_b\simeq0_{\mathcal A}$), plus la contrainte supplémentaire $\mathbb D$-cartésienne, sur le carré cherché, détermine celui-ci à isom unique près en termes de son squelette : il y a une substantiation non seulement ``canonique'' de celui-ci, mais \unsure{$=$.} unique à isom unique près, \deleted{\unsure{en termes}} à cause de la contrainte $\mathbb D$-cartésienne.

Je voudrais commencer par regarder le cas plus simple encore où

\diagram{\begin{tikzcd}[ampersand replacement=\&]
I=\mathbb{\Delta}_1:\quad 0\arrow[r]\&1
\end{tikzcd}\,,}

donc \deleted{\unsure{les diagrammes}} les squelettes envisagés sont simplement \unsure{des|les} flèches

\diagram{\begin{tikzcd}[ampersand replacement=\&] F_1\arrow[r]\&F_0\end{tikzcd}\quad dans $\mathcal A$.}

La philosophie nous oblige, pour les flèches les plus fondamentales dans la \unsure{théorie} du dérivateur $\mathbb D$ (flèches dans $\mathcal A$, ou par extension dans une $\mathbb D(I)\simeq\mathcal A_{\mathbb D_I}$, cat. fondamentale du \deleted{\unsure{$\mathcal A$}} dérivateur \unsure{induit} $\mathbb D_I$), de se demander s'il n'y a pas
